\documentclass[12pt]{article}
\usepackage[utf8]{inputenc}
\usepackage[T1]{fontenc}
\usepackage{hyperref}
\title{Western Tlacolula Valley Zapotec Language Resources for K-12 Curriculum Development}
\author{Compiled by Claude Sonnet 5.5 (claude-sonnet-5-5) for Andr\'as Kornai}
\date{}
\begin{document}
\maketitle

\section{Introduction}
Western Tlacolula Valley Zapotec (ISO 639-3 zab; Glottocode sanj1284) is an Otomanguean language of Oaxaca, Mexico, with about 29,000 speakers. It is catalogued with San Juan Guelav\'ia. My searches found few dedicated resources and no school materials.

\section{Primary Educational Resources}
None identified in my searches.

\section{Linguistic and Reference Materials}
\begin{itemize}
\item Joseph and Judith Bradley Whitecotton, \emph{Vocabulario zapoteco-castellano} (1993, Vanderbilt University Publications in Anthropology), catalogued under zab. Glottolog record: \url{https://glottolog.org/resource/reference/id/305642}
\item A 1984 collection of anecdotes told by Pedro Hern\'andez in Guelav\'ia Zapotec with Spanish translation (ILV, Mexico City). The search result for it was a repository file: \url{https://repositorio.colmex.mx/downloads/nz8063419?locale=en}
\item Project MUSE article, found in search for Guelav\'ia and Valley Zapotec (I did not confirm its content): \url{https://muse.jhu.edu/article/19092/summary}
\end{itemize}

\section{Religious and Community Resources}
None identified in my searches.

\section{Digital and Audio Resources}
\begin{itemize}
\item Dictionary of the Zapotec Language by \'Oscar M\'endez Espinoza (five volumes, many communities, not confirmed to include Guel\'avia), hosted by AILLA: \url{https://islandora-ailla.lib.utexas.edu/islandora/object/ailla%3A279600}
\item Ticha, Haverford College, a digital resource on Valley Zapotec (mainly colonial texts, a different variety): \url{https://ds-omeka.haverford.edu/ticha/en/index.html}
\end{itemize}

\section{Curriculum Development Context}
With about 29,000 speakers the language has a sizeable community, but published teaching material is scarce. A curriculum would build on the Whitecotton vocabulary, the ILV texts and new work with speakers. It would also need a decision on orthography.

\end{document}